\documentclass[11pt]{article}

\usepackage[margin=1in]{geometry}
\usepackage{amsmath,amssymb,amsthm,mathtools,mathrsfs,booktabs,tabularx,longtable}
\usepackage{microtype,xurl}
\usepackage{xcolor}
\usepackage[colorlinks=true,linkcolor=blue!55!black,
            citecolor=blue!55!black,urlcolor=blue!55!black]{hyperref}

\newcommand{\dd}{\mathrm d}
\newcommand{\R}{\mathbb R}
\newcommand{\Diff}{\operatorname{Diff}}
\newcommand{\im}{\operatorname{im}}
\newcommand{\cert}[1]{\nolinkurl{#1}}
\newcolumntype{Y}{>{\raggedright\arraybackslash}X}
\newenvironment{resultcard}[1]
 {\par\medskip\noindent\hrule\smallskip\noindent\textbf{#1}\par\smallskip}
 {\par\smallskip\hrule\medskip}

\title{What Survives the Ghost Test?\\
\large A Four-Level Phase-1 Classification in Pure-Weyl Gravity}
\author{GPT-5.6.sol\thanks{OpenAI model and principal programme author.}
\and Claude Fable 5\thanks{Anthropic model and computational coauthor. The
project was commissioned and coordinated by Asger Alstrup Palm
(\texttt{asger@area9.dk}), who initiated the questions, exercised editorial
control, and serves as corresponding human contact.}}
\date{Phase-1 synthesis, 22 July 2026}

\begin{document}
\maketitle

\begin{abstract}
Pure Weyl gravity is a fourth-order conformal metric theory.  Its linearized
equations contain more solutions than Einstein gravity, and standard
propagator decompositions assign indefinite signs to some of them.  Those
signs diagnose a problem, but do not by themselves determine a physical state
space.  We separate four logically different objects: local solutions;
gauge-, charge-, and boundary-reduced classical directions; nonlinearly
continuable directions; and quantum states or particles.

On compact Weyl--Maxwell laboratories, additional axial and polar directions
survive local gauge reduction and have nonzero induced Lee--Wald pairings.
Their classical signs are not organized by the simple rule ``Einstein
positive, additional negative.''  At second order, finite-harmonic
exponential-polynomial continuation is equivalent to five quadratic
stabilizer-charge conditions, whereas bounded continuation obeys further
polynomial-growth and characteristic-shell resonance constraints.  On
Schwarzschild, horizon regularity alone admits additional curvature
solutions, while finite symplectic norm excludes the additional axial and
polar solutions at the certified $(\ell,\omega)=(2,3/5)$ fixture.  In the
quantum analysis, strict fixed-field-content pure Weyl gravity has a nonzero
local Euclidean one-loop BV anomaly.  A formal compensator makes the anomaly
exact, but changes the theory and does not supply a Lorentzian state space.

We also analyze a changed two-phase counterflow gravity--clock theory.  It has
an exact 70-component causal BV parent and removes a dressed-trace
obstruction.  Charge reduction reveals an exact tradeoff: making time
translation null removes the relative clock, while retaining the clock makes
time translation charged.  More decisively, the complete
$j=\tfrac12$ reduced block contains a Hamiltonian--Hopf quartet throughout
the connected trace-healthy same-field stationary family.  Phase 1 therefore
ends with a classification rather than a selected theory: pole-level signs
are too coarse, but the more complete reductions do not generically rescue
the candidates tested here.  No positive full-BV state, scattering theory,
or unitarity theorem is claimed.
\end{abstract}

\tableofcontents

\section{The question and the four levels}

The metric action of strict pure Weyl gravity is
\begin{equation}
 S_{\rm W}[g]
 =\alpha_{\rm W}\int_M \sqrt{-g}\,
 C_{\mu\nu\rho\sigma}C^{\mu\nu\rho\sigma}\,\dd^4x,
 \qquad B_{\mu\nu}=0,
 \label{eq:weyl-action}
\end{equation}
where $B_{\mu\nu}$ is the Bach tensor.  Around a fixed background, the
linearized Bach equation is fourth order.  A formal factorization, a partial
fraction decomposition, or a negative residue can reveal an indefinite
direction.  None of these alone says whether the direction survives gauge
reduction, global constraints, physical boundary conditions, nonlinear
integrability, or the construction of a quantum state.

This paper organizes the first phase of a computer-assisted research
programme by the following hierarchy:
\begin{equation}
\boxed{
\begin{array}{cl}
\text{Level 1:}&\text{local solutions of the linear field equations},\\
\text{Level 2:}&\text{reduced classical directions after gauge, charge,
and boundary conditions},\\
\text{Level 3:}&\text{directions tangent to nonlinear solutions in a stated
correction class},\\
\text{Level 4:}&\text{quantum states, particles, and a positive physical
probability interpretation}.
\end{array}}
\label{eq:four-levels}
\end{equation}
The arrows between these levels are mathematical problems.  They are not
terminological promotions.  A nonzero classical symplectic pairing at Level
2 is not a positive quantum norm at Level 4.  A secular second-order
primitive at Level 3 is not nonlinear stability.  A local Euclidean anomaly
calculation is not a Lorentzian scattering theorem.

The phase-one question is therefore not simply ``does the propagator contain
a ghost?''  It is:
\begin{quote}
Which additional fourth-order directions survive each declared reduction,
which fail nonlinear or boundary tests, and what remains undecided before a
particle interpretation can even be posed?
\end{quote}

The answer is mixed.  Some additional directions survive local gauge and
presymplectic reduction.  Compact nonlinear constraints remove some of them,
while balanced combinations can continue formally with secular terms.  A
specified Schwarzschild finite-norm condition removes them at one two-parity
fixture.  The strict quantum gauge theory is locally anomalous.  Finally, a
changed theory that passes a difficult causal construction fails a physical
spectral test.  Different candidates fail at different levels.

\section{Theories, conventions, and scope}

We use Lorentzian signature $(-,+,+,+)$ and
\(
[\nabla_\mu,\nabla_\nu]v^\rho
=R^\rho{}_{\sigma\mu\nu}v^\sigma
\).
Classical fields are real; complex Fourier modes are bookkeeping devices and
physical pairings use the corresponding real or conjugate-frequency form.
The connected local gauge group of the strict theory is
\begin{equation}
 \Diff_0(M)\ltimes C^\infty(M,\R)_{\rm Weyl},
 \qquad g_{\mu\nu}\mapsto e^{2\sigma}g_{\mu\nu}.
\end{equation}
Unless a theorem card says otherwise, no boundary counterterm or topological
density is included in the action.  Boundary and support conditions are part
of the theorem, not implicit conventions.

Two comparison theories appear.  On the compact product we use
\begin{align}
 S_{\rm WM}[g,A]
 &=S_{\rm W}[g]-\frac14\int_M\sqrt{-g}\,F_{\mu\nu}F^{\mu\nu}\,\dd^4x,
 \\
 S_{\rm EM}[g,A]
 &=\frac1{2\kappa}\int_M\sqrt{-g}\,(R-2\Lambda)\,\dd^4x
   -\frac14\int_M\sqrt{-g}\,F_{\mu\nu}F^{\mu\nu}\,\dd^4x.
\end{align}
The magnetic bundle is fixed in that comparison.  The Einstein--Maxwell
theory is a source/comparator; it is not silently identified with a subsector
having the same symplectic structure.

The quantum repair uses a formal Wess--Zumino field $\tau$ with
\(
s\tau=\omega
\), where $\omega$ is the Weyl ghost, and the dressed metric
\(
\widehat g=e^{-2\tau}g
\).
This is an enlarged BV theory.  It is not strict pure Weyl gravity in another
gauge.

The later counterflow experiment changes the classical action as well.  Its
defining matter mechanism is two positive phase fields with an auxiliary
diagonal $U(1)$ connection,
\begin{equation}
 S_{\rm phase}
 =-\frac12\int_M\sqrt{-\widehat g}\,
 \left[f_1^2(\dd\theta_1-A)^2+f_2^2(\dd\theta_2-A)^2\right].
 \label{eq:counterflow-action}
\end{equation}
Writing
\begin{equation}
 \chi=\frac{f_1^2\theta_1+f_2^2\theta_2}{f_1^2+f_2^2},
 \qquad \psi=\theta_1-\theta_2,
 \qquad
 \mu^2=\frac{f_1^2f_2^2}{f_1^2+f_2^2},
\end{equation}
separates a diagonal gauge phase from a gauge-invariant relative phase with
positive bare kinetic coefficient $\mu^2$.  The certified model also contains
the action-derived repaired compensator/scale sector recorded by its frozen
70-component Hessian.  Results for this changed theory are never imported
back into the strict theory.

\subsection{The scope tuple}

Every result is typed by
\begin{equation}
 \boxed{\mathfrak S=
 (\text{theory},\text{background},\text{generator},
  \text{phase space},\text{boundary/support class},\text{claim state}).}
 \label{eq:scope-tuple}
\end{equation}
Two statements with different entries in $\mathfrak S$ are not combined
without a constructed map.  This matters because the programme uses several
controlled laboratories:
\begin{center}
\begin{tabularx}{\textwidth}{@{}lY@{}}
\toprule
Laboratory & Question isolated there \\
\midrule
$\R\times S^3$ conformal cylinder & Complete free causal complex and a
selected residual zero-charge reduction \\
Biaxial Berger spacetime & Clocks, relational observables, interactions, and
the counterflow successor \\
$\R\times S^1\times S^2$ magnetic product & Einstein/additional branches,
Lee--Wald pairing, Taub constraints, and resonances \\
Nariai & Transfer of the causal construction to a second curved background \\
Schwarzschild and static spherical metrics & Horizons, finite-flux
admissibility, and black-hole perturbations \\
Euclidean $S^4$ and regular local charts & One-loop determinants and local BV
anomaly classes \\
\bottomrule
\end{tabularx}
\end{center}
These are related tests of a common architecture, not successive
approximations to one already unified universe.

\section{What ``ghost'' can mean}

The literature uses the word \emph{ghost} for inequivalent failures.  The
distinctions are especially important in a gauge theory with a fourth-order
equation.

\begin{longtable}{@{}p{0.22\textwidth}p{0.39\textwidth}p{0.31\textwidth}@{}}
\toprule
Diagnosis & Mathematical question & Phase-1 disposition \\
\midrule
\endhead
Ostrogradsky direction & Is a nondegenerate higher-derivative Hamiltonian
unbounded before constraints? & Crosswalks are computed on selected blocks;
this is not used as a complete physical quotient. \\
Negative pole residue & Does a gauge-fixed propagator decompose with a
negative coefficient? & Treated as an initial warning, not a final state
classification. \\
Negative classical energy or flux & What is the sign of the descended
Lee--Wald/canonical form on a declared real quotient? & Computed for selected
compact and black-hole sectors.  Some signs contradict a naive
Einstein/additional split. \\
Radical or pure gauge & Does the induced presymplectic form vanish on the
direction? & Additional compact-product directions are nonradical on the
declared pre-residual quotient. \\
Nonlinear obstruction & Is a linear direction tangent to a second-order
solution in a stated function class? & Classified completely on a finite
harmonic exponential-polynomial space; bounded continuation is stricter. \\
Negative quantum norm & Is there a BRST-compatible positive state on the
full physical algebra? & Not established.  Reduced indefinite two-point
functions are not a full-BV state. \\
Failure of unitarity & Does an asymptotic state space satisfy the optical
theorem? & Not tested: no complete scattering or particle theory exists. \\
Alternative pole prescription & Is the theory redefined by a PT-symmetric,
Dirac--Pauli, or fakeon prescription? & Not adopted.  These are neighboring
programmes with different physical rules, not consequences of our classical
reductions. \\
\bottomrule
\end{longtable}

PT-symmetric treatments of the Pais--Uhlenbeck system, Dirac--Pauli
formulations of quadratic gravity, and fakeon prescriptions all demonstrate
that ``negative residue implies negative-probability asymptotic particle'' is
not the only proposed logical route
\cite{BenderMannheim,Salvio,Anselmi}.  They also demonstrate why the chosen
state space and prescription must be declared.  Phase 1 follows the ordinary
Lorentzian BV/Lee--Wald route and leaves those alternatives as comparisons.

\section{Which additional linear directions survive reduction?}

The compact magnetic product is the cleanest place to compare the
Einstein--Maxwell and Weyl--Maxwell linear solution spaces.  The computation
uses both parities, generic harmonics, exceptional modes, and global blocks.

\begin{resultcard}{Result card A: compact Einstein/additional decomposition}
\textbf{Theory.} Einstein--Maxwell source and Weyl--Maxwell target.\par
\textbf{Background.} The magnetically supported compactified
Plebanski--Hacyan product $\R\times S^1\times S^2$, in a fixed magnetic
bundle.\par
\textbf{Domain.} Parity-complete generic linear modes, together with the
declared exceptional and global blocks.\par
\textbf{Gauge and quotient.} Local gauge quotient and maximal pre-residual
$H^0$ exact sequence; final stabilizer reduction is not included.\par
\textbf{Boundary condition.} Spatial periodicity; no asymptotic scattering
condition.\par
\textbf{Statement.} The Einstein--Maxwell solutions inject into the
Weyl--Maxwell solutions.  The cofiber contains additional axial and polar
blocks.  Each additional generic block is nonradical with Lee--Wald inertia
$(2,0)$, while the target Einstein image has inertia $(1,1)$ and the complete
generic target has inertia $(3,1)$.  The source Einstein--Maxwell form has
inertia $(2,0)$, so inclusion does not define a cyclic or symplectic
equivalence.\par
\textbf{Not claimed.} No positive energy, particle norm, final residual
quotient, causal boundary admissibility, or nonlinear closure follows.
\end{resultcard}

The important conclusion is not that the additional block is ``healthy''
because its displayed inertia is positive.  The conclusion is that it is not
an algebraic duplicate or a radical direction on this quotient, and that the
standard pairing does not respect the naive expectation that the Einstein
image should inherit the ordinary Einstein sign while the complement should
carry the bad sign.  The action matters: the Einstein--Hilbert and Weyl
Lee--Wald forms are different even on the same metric perturbation.

This exact-sequence formulation also prevents an unjustified direct-sum
claim.  A lift of an additional curvature representative may be shifted by
an Einstein solution.  The invariant information is the injection, cofiber,
induced ranks, radicals, and extension class.  A canonical branch split
requires more structure than composition of differential operators.

On the conformal cylinder a different reduction occurs.  In a selected
derived zero-charge sector, all fifteen residual conformal generators are
implemented by a BFV/Koszul receiver, and the vacuum and selected residual
one-particle cohomologies are acyclic.  The first surviving $H^4$ classes are
deformation or vertex classes, not one-particle gravitons.  This is a theorem
about that selected compact sector.  It is not a universal removal of
radiative gravitons and it supplies no asymptotic particle interpretation.

The two laboratories therefore teach complementary lessons: additional
directions can survive a compact local-gauge quotient, while a stronger
zero-charge residual reduction can remove selected one-particle
cohomology.  Neither statement transfers to the other without a map.

\section{Which directions survive nonlinear integrability?}

Linearization stability on compact Cauchy surfaces has a classical history
from Brill--Deser through Fischer--Marsden, Moncrief, and the
Arms--Marsden--Moncrief momentum-map normal form
\cite{BrillDeser,FischerMarsden,Moncrief,AMM}.  The general idea that a
symmetric solution has a quadratic momentum-map cone is therefore not new.
The contribution here is an exhaustive model-specific calculation for a
fourth-order Weyl--Maxwell mode space and a separation of formal secular
solvability from bounded solvability.

Let $L$ be the linearized equation and $D^2E[u,u]$ the quadratic source.  Two
correction spaces are kept distinct:
\begin{align}
 \mathscr X_{\rm ep}
 &=\left\{\text{finite sums of }t^p e^{-i\Omega t}
    \text{ with spatially periodic finite harmonic support}\right\},\\
 \mathscr X_{\rm qp}
 &=\left\{\text{finite bounded quasiperiodic sums }
    e^{-i\Omega t}\right\}.
\end{align}
The first allows secular polynomials; the second does not.

\begin{resultcard}{Result card B: finite-harmonic second-order cone}
\textbf{Theory.} Compact Weyl--Maxwell theory, with the Einstein--Maxwell
comparison used to label branches.\par
\textbf{Background.} $\R\times S^1\times S^2$ with the fixed magnetic
bundle.\par
\textbf{Domain.} The complete declared finite harmonic space: generic
Einstein and additional primaries, both parities, exceptional dipoles,
twists, homogeneous generalized modes, electric flux, and Wilson directions.\par
\textbf{Gauge and quotient.} The five stabilizer covectors generated by time
translation $H$, circle translation $P_x$, and rotations $J_1,J_2,J_3$.\par
\textbf{Correction class.} Finite exponential-polynomial corrections
$\mathscr X_{\rm ep}$.\par
\textbf{Statement.}
\[
 \mathcal Z^{\rm ep}_2
 =\{u:\mu_H(u)=\mu_{P_x}(u)=
          \mu_{J_1}(u)=\mu_{J_2}(u)=\mu_{J_3}(u)=0\}.
\]
The five moment maps are necessary and sufficient in this correction class.\par
\textbf{Not claimed.} This is not a theorem for all smooth or Sobolev fields,
retarded solutions, bounded dynamics, all-order integration, or stability.
\end{resultcard}

The sufficiency statement follows from blockwise constant-coefficient
surjectivity: nonzero invariant factors are surjective on finite
exponential-polynomial modules, including at characteristic roots after the
appropriate polynomial degree is allowed.  The only zero invariant factors
are the five stabilizer covectors.  This is the explicit exhaustive content
beyond the general momentum-map normal form.

Bounded continuation has a larger obstruction ledger.  If the quadratic
source is expanded by frequency shell and polynomial degree, then a bounded
primitive exists only when three types of functional vanish:
\begin{equation}
 \{\mu_X\}
 \ \cup\ 
 \{\mathcal P_{j,r}\}
 \ \cup\ 
 \{\mathcal R_{j,a}\}.
 \label{eq:bounded-ledger}
\end{equation}
Here $\mu_X$ are the global zero-block Taub constraints,
$\mathcal P_{j,r}$ are coefficients that would require positive powers of
$t$, and $\mathcal R_{j,a}$ are adjoint-kernel pairings on nonzero
characteristic shells.  Their common nonlinear zero locus is not fully
classified for arbitrary finite support.  It is classified on several
important subspaces.  In particular, on the complete global plus
$(\ell=2,k=0)$ additional-primary space, the bounded cone reduces to the
standard global generalized-zero cone: no nonzero additional-primary
amplitude survives.  At other momenta, balanced sheets can remain.

The distinction is physical.  A resonant source can have a formal primitive
$t e^{-i\omega t}$ and hence define a second-order jet, while failing every
bounded quasiperiodic test.  The programme therefore does not report
``second-order solvable'' without naming the correction class.

The same calculation also blocks a common branchwise simplification.
Balanced Einstein--additional configurations can have zero total moment map
while their projected Einstein and additional pieces carry equal and
opposite nonzero charges.  Consequently the ambient exact sequence does not
descend to separately neutral branch fibers.  Nonlinear reduction is
intrinsically mixed.

\section{Which boundary conditions retain the extra solutions?}

Compact charge constraints and black-hole endpoint conditions test different
questions.  The Schwarzschild analysis uses the Ricci perturbation
$\psi_{\mu\nu}=\delta R_{\mu\nu}[h]$ as a second-order curvature variable in
the fourth-order Bach system.  On a Ricci-flat background the universal
linear identity is
\begin{equation}
 \delta B_{\mu\nu}
 =\frac12\Box\delta R_{\mu\nu}
  +C_{\mu\rho\nu\sigma}\delta R^{\rho\sigma}
  -\frac16\nabla_\mu\nabla_\nu\delta R
  -\frac1{12}g_{\mu\nu}\Box\delta R.
 \label{eq:ricci-bach}
\end{equation}
The Einstein image lies in $\ker\delta R$; the additional curvature content
is visible in $\psi$.

Horizon analyticity does not force $\psi$ to vanish.  Regular ingoing
additional curvature solutions exist in the certified real-frequency axial
analysis.  This already shows that future-horizon regularity alone does not
select the Einstein image.  It does not show that every local differential
initial or boundary condition fails: imposing $\psi$ and its normal
derivative to vanish on a Cauchy surface is a local restriction that selects
the Einstein image at linear order.

The outer boundary gives a sharper fixture-level result.

\begin{resultcard}{Result card C: finite-norm selection at a Schwarzschild fixture}
\textbf{Theory.} Strict pure $C^2$ gravity.\par
\textbf{Background.} Schwarzschild exterior.\par
\textbf{Domain.} Axial and polar $\ell=2$ reconstructed metric solutions at
real frequency $\omega=3/5$, with the declared Lee--Wald flux/norm
normalization.\par
\textbf{Gauge and quotient.} Regge--Wheeler/Zerilli-type parity reduction and
the action-derived pure-Weyl Lee--Wald current.\par
\textbf{Boundary condition.} Future-horizon ingoing regularity followed by
finite symplectic norm in the certified asymptotic radiation class.\par
\textbf{Statement.} At this two-parity fixture the Einstein solutions have
finite norm, whereas the reconstructed additional axial and polar solutions
have divergent norm.  Finite norm therefore selects the Einstein sector at
the stated $(\ell,\omega)$.\par
\textbf{Not claimed.} No all-frequency theorem, full exterior initial-boundary
well-posedness, quasinormal spectrum, ringdown, stability, particle
interpretation, or universal causal-truncation theorem follows.
\end{resultcard}

This result is deliberately narrower than the proposition that ``ordinary
boundary conditions always remove the additional branch.''  The latter
would require a complete asymptotically flat Bach phase space, metric
reconstruction and Jordan analysis at all relevant frequencies, finite flux,
and a well-posed exterior problem.  Existing AdS/Euclidean Einstein-selection
arguments operate in different boundary classes
\cite{Maldacena,LuPangPope}.  The present contribution is a Lorentzian
Schwarzschild horizon/finite-norm comparison at an exact declared fixture.

\section{Can a positive quantum state be constructed?}

There are three separate quantum objects in the programme: local BV anomaly
cohomology, Euclidean one-loop coefficients, and reduced Lorentzian
two-point functions.  They are not interchangeable.

In four dimensions a local Weyl breaking at ghost number one and form degree
four has essential type-B and type-A representatives of the form
\begin{equation}
 \mathcal A
 =\int_M\sqrt g\,\omega
 \bigl(c_{\rm W}C_{\mu\nu\rho\sigma}C^{\mu\nu\rho\sigma}
       -a_{\rm W}E_4+b_{\rm W}\Box R\bigr)\,\dd^4x .
 \label{eq:anomaly}
\end{equation}
The $\Box R$ coefficient is scheme dependent and may be shifted by a local
counterterm.  The nonzero $C^2$ and Euler representatives are the essential
local classes.  A one-loop divergence, beta function, trace anomaly, and BV
master-equation breaking are related only after the gauge-fixed determinant,
normalization, descent, and quotient maps have been supplied.

\begin{resultcard}{Result card D: strict and compensated local quantum theories}
\textbf{Strict theory.} Fixed-field-content $\Diff\ltimes$Weyl pure $C^2$
gravity.\par
\textbf{Domain.} Complete declared gauge-fixed local BV quotient on a regular
Bach-locus chart, with Euclidean one-loop coefficient input.\par
\textbf{Statement.} The regulated breaking has nonzero components along the
essential ghost-number-one, form-degree-four Weyl classes.  The strict local
Euclidean one-loop QME is obstructed.\par
\textbf{Extended theory.} In the formal $\tau$-adic Wess--Zumino BV algebra
with $s\tau=\omega$, the same cocycle is exact and a local counterterm restores
the Euclidean QME at one loop.\par
\textbf{Not claimed.} The extension changes the BV theory.  It gives no
unconditional all-loop theorem, regulator/QAP for an arbitrary changed
action, Lorentzian QME, Hadamard state, positive physical Hilbert space,
particles, scattering, or unitarity.
\end{resultcard}

The compensator statement is an exact BV realization of the familiar
dilaton/Wess--Zumino mechanism, whose broader structure is known
\cite{BaumeKerenZur}.  Its value is not a claim that a compensator is a free
repair.  Indeed, the passive compensator leaves a gauge-invariant dressed
trace in the classical causal complex.  Repairing that obstruction led to
the changed counterflow experiment in Section~\ref{sec:counterflow}.

Separately, the programme constructs candidate Hadamard two-point functions
on a reduced free Einstein/additional/longitudinal mode space.  The candidate
pairing has signs $(+,-,-)$ on that selected space.  This is useful evidence
that an indefinite direction survives that reduction.  It is not a
BRST-compatible Hadamard covariance for the full metric BV complex and not a
one-particle Hilbert space.  No positive full-BV state has been constructed.

The immediate quantum literature includes conformal-gauge fixing and its
additional ghosts, local BRST cohomology, and several alternative
quantization prescriptions\cite{Rachwal,BarnichBrandtHenneaux}.  Phase 1
neither reproduces all of that literature nor supersedes it.  It contributes
a content-addressed strict/extended local-BV calculation and records the
precise point at which Lorentzian state construction remains absent.

\section{What the counterflow candidate taught us}
\label{sec:counterflow}

The strict anomaly and passive-compensator causal obstruction motivated a
small changed-theory ladder.  Passive trace repairs and a one-phase active
clock failed combined causal, sign, and charge tests.  Two phases with an
auxiliary diagonal $U(1)$ allow neutral counterflow:
\begin{equation}
 \dot\chi=0,
 \qquad \dot\psi=\Omega,
 \qquad Q_{\rm diag}=0,
 \qquad E_{\rm rel}=\frac12\mu^2\Omega^2>0.
\end{equation}
The relative phase can move while the diagonal compact gauge charge
vanishes.  Because its stress tensor is built from the dressed metric, it is
quadratically visible in the dressed-trace direction.

\begin{resultcard}{Result card E: causal construction and physical nonselection}
\textbf{Theory.} The action-derived two-phase counterflow theory with
auxiliary diagonal $U(1)$ and repaired compensator sector.\par
\textbf{Background.} The selected positive biaxial Berger solution and its
connected trace-healthy same-field stationary retuning family.\par
\textbf{Domain.} A 70-component cyclic BV parent, decomposing at the selected
point into a 54-component causal parent and a 16-component contractible
diagonal-$U(1)$ sector; physical spectral reduction uses complete
$SU(2)_L\times U(1)_R$ isotypical blocks.\par
\textbf{Gauge and charge.} Diagonal gauge charge, relative-phase charge
$Q_{\rm rel}$, time translation $D$, and the background-preserving generator
$K=D-\Omega R_{\rm rel}$ are kept distinct.\par
\textbf{Causal statement.} Exact retarded and advanced homotopies, nilpotency,
cyclicity, and support identities hold at the selected point, and the original
dressed-trace cohomology class is removed.\par
\textbf{Physical statement.} The complete both-$k$, $j=\tfrac12$ reduced
block contains a Hamiltonian--Hopf quartet.  The obstructing factor is
\[
 (40z^4+773z^2+3748)^2,
 \qquad 773^2-4(40)(3748)=-2151<0.
\]
The quartet persists throughout the connected trace-healthy same-field
stationary family.  Therefore no member of that family is selected as a
robust linearly stable Phase-2 candidate.\par
\textbf{Not claimed.} This is not an all-isotype spectrum, a familywide
Green-homotopy theorem, a full-PDE stability theorem, or a no-go over other
field contents and gauge architectures.
\end{resultcard}

The rejection criterion was fixed before the computation: a Phase-2
candidate had to be linearly stable in every physical sector.  One exact
unstable physical isotypical block is sufficient to reject it.  It is not
necessary to finish every higher-$j$ block to establish this nonselection.
Conversely, the finite block does not by itself describe the nonlinear
evolution of arbitrary PDE data.

The quartet is a spectral Hamiltonian instability associated with a
Krein-sign collision, not merely a coordinate drift.  The global clock block
has a different status.  On the unrestricted phase space,
\begin{equation}
 \Omega_{\rm rel}=\dd Q_{\rm rel}\wedge\dd\psi,
 \qquad
 H_{\rm rel}=\frac{Q_{\rm rel}^2}{2I}+E_{\rm geometry},
\end{equation}
so its secular dephasing is tangent to an action--angle family and passes an
orbital, frequency-modulated interpretation.  It does not cure the
$j=\tfrac12$ quartet.

The charge analysis gives an independent theorem.  For phase clocks with
linear charge constraint $CQ=q_0$ and clock velocity vector $v$,
\begin{equation}
 D\text{ is null after reduction}
 \quad\Longleftrightarrow\quad
 v\in\im C^T,
\end{equation}
while physical clocks are classes of $v$ modulo $\im C^T$.  In the certified
model, fixing and quotienting $Q_{\rm rel}$ makes $D$ null but contracts the
relative-clock Darboux pair.  Keeping the clock leaves
\(
H_D=\Omega Q_{\rm rel}
\)
and raw $D$ charged; $K=D-\Omega R_{\rm rel}$ preserves the background.  No
single declared charge fibre supplies both the proposed physical clock and a
gauge time translation.

The counterflow result is therefore not a failed piece of engineering.  It
separates three notions that are often conflated:
\begin{equation}
 \begin{aligned}
 \text{causal completeness}
 &\quad\not\Rightarrow\quad
 \text{positive reduced dynamics},\\
 \text{positive reduced dynamics}
 &\quad\not\Rightarrow\quad
 \text{acceptable clock/charge interpretation}.
 \end{aligned}
\end{equation}
The causal theorem remains valid even though the candidate is rejected by a
later physical gate.

\section{Integrated Phase-1 verdict}

The main findings can now be placed on the four levels of
Eq.~\eqref{eq:four-levels}.

\begin{center}
\begin{tabularx}{\textwidth}{@{}p{0.16\textwidth}YY@{}}
\toprule
Level & Established in a declared scope & Still absent \\
\midrule
Local linear systems & Complete causal BV complexes on selected compact
backgrounds; exact additional compact-product and Schwarzschild curvature
solutions & A single causal theory covering every laboratory and boundary
class \\
Reduced classical directions & Additional axial/polar compact-product
directions are nonradical; counterflow causal parent and exact charge
classification; finite-norm Schwarzschild fixture & Universal asymptotic
phase space, positive energy theorem, and acceptable counterflow spectrum \\
Nonlinear continuation & Five-charge necessary-and-sufficient finite
exponential-polynomial cone; bounded resonance ledger and classified
subcones & Full bounded finite-support zero locus, retarded nonlinear
existence, all-order integration, and stability \\
Quantum states/particles & Strict local one-loop BV obstruction; formal
changed-theory Wess--Zumino resolution; reduced indefinite two-point
functions & Lorentzian QME, full-BV Hadamard state, positive physical Hilbert
space, particles, scattering, optical theorem, and unitarity \\
\bottomrule
\end{tabularx}
\end{center}

This supports the following central claim:
\begin{quote}
\textbf{A pole-level diagnosis is too coarse, but the more complete reduction
does not generically rescue the theory.}  Some additional classical
directions survive gauge and presymplectic reduction.  Nonlinear moment maps,
resonances, boundary norms, local anomalies, and reduced spectral tests then
eliminate different candidates for different reasons.  The first changed
theory to achieve a complete causal BV parent is rejected throughout its
connected trace-healthy same-field family by a physical
$j=\tfrac12$ Hamiltonian--Hopf instability.
\end{quote}

This is a classification ending, not a viable-theory claim and not a
universal no-go.  A different field content, a different gauge architecture,
or a different admissible state prescription could define a new Phase 2.
Phase 1 does not select one.

\section{Relation to prior work}

The programme combines several established structures.  Its claims should be
read as explicit constructions and exhaustive calculations in specified
models, not as inventions of the surrounding general theories.

\paragraph{Causal gauge complexes.}
Green-hyperbolic complexes, retarded/advanced Green homotopies, and their
induced Poisson structures have a general Lorentzian theory
\cite{BeniniMusanteSchenkel}.  Conformal deformation and Yang--Mills detour
complexes place the Bach operator in a longer conformal-geometric lineage
\cite{BransonGover,GoverSombergSoucek}.  The contribution here is the explicit
action-derived conformal-gravity BV construction, its cyclic pairings and
support identities, and transfer across the declared backgrounds.

\paragraph{Linearization stability.}
The relation between Killing symmetries, quadratic Taub constraints, and
singular momentum-map zero sets is classical
\cite{BrillDeser,FischerMarsden,Moncrief,AMM}.  The finite-harmonic theorem is
an exhaustive Weyl--Maxwell realization: exactly five covectors exhaust the
exponential-polynomial cokernel, exceptional and generalized modes are
included, and bounded resonance conditions are separated from the formal
cone.

\paragraph{Einstein selection and black holes.}
Boundary-condition selection of Einstein solutions in conformal gravity is
well known in Euclidean and asymptotically (A)dS settings
\cite{Maldacena,LuPangPope,DeserJoungWaldron}.  Prior Weyl-black-hole
quasinormal analyses often isolate a second-order conformally transported
sector rather than the full fourth-order Bach solution space
\cite{MomenniaHendi}.  Our fixture-level finite-norm result addresses the
full additional reconstruction in a one-ended Lorentzian Schwarzschild
problem, but does not yet provide a general asymptotically flat Bach phase
space.

\paragraph{Quantum conformal gravity.}
One-loop divergences, conformal anomalies, BRST cohomology, and conformal
gauge fixing have extensive prior literatures
\cite{FradkinTseytlin,BarnichBrandtHenneaux,Rachwal}.  The phase-one quantum
claim is narrower: an exact local-BV quotient and coefficient ledger for the
declared strict theory, followed by a formal $\tau$-adic Wess--Zumino
trivialization in a changed theory.  It does not settle the competing
PT-symmetric, Dirac--Pauli, fakeon, or conventional positive-state proposals.

\paragraph{Relational clocks.}
The construction of complete or partial observables from internal clocks is
standard constrained-system territory\cite{Dittrich}.  The counterflow
experiment adds a typed distinction between diagonal gauge charge, relative
clock charge, time translation $D$, and the background-preserving combination
$K$.  No operational counterflow frequency ratio survived Phase 1, so the
result is a charge/clock classification rather than an observational
prediction.

\section{Limitations and the next scientific boundary}

The classification leaves several sharp, finite questions.  The most
important are not ``finish quantum gravity.''  They are:
\begin{enumerate}
\item construct a complete asymptotically flat Bach phase space and determine
which additional black-hole solutions have finite flux for general
frequency and angular momentum;
\item determine the common zero locus of all bounded nonlinear obstruction
functionals on the full declared finite-support space;
\item decide whether the retained mixed interaction defines a nonzero class
in the complete bounded cyclic deformation complex (its physical action is
trivialized through summed input order two, so the earlier representative is
not yet an invariant physical interaction);
\item construct a BRST-compatible Hadamard covariance for a full Lorentzian
BV complex, or prove a scoped obstruction;
\item if a new changed theory is proposed, require an open stationary locus
with causal completeness, positive reduced dynamics, coherent charge
interpretation, a preserved nonlinear sector, and at least one operational
observable before calling it a candidate.
\end{enumerate}

These questions define a possible Phase 2.  They are not conditions for
closing Phase 1.  The stopping rule was outcome-independent: Phase 1 ended
when the programme could give an exact viability classification of the first
causally complete changed-theory candidate and reconcile it with the compact
linear, nonlinear, black-hole, observer, and quantum results.  That record is
now frozen.

\section{Computational accountability}

The project treats computer algebra as an experimental instrument whose
outputs require independent checks.  Material results are represented by a
producer, a machine-readable certificate, a verifier that does not merely
rerun the producer, adversarial mutation tests where practical, a tiered test
receipt, and a human-readable scope report.  Exact algebraic ranks, Smith
forms, cohomology witnesses, and characteristic factors use rational or
algebraic arithmetic; numerical evidence is typed separately.

The authoritative Phase-1 record is
\begin{center}
\cert{PURE_WEYL_PROGRAMME_PHASE1_CLASSIFICATION_ENDING_V1},
\end{center}
stored in
\begin{center}
\cert{reports/phase1-closure-claims-ledger-2026-07-22.json}.
\end{center}
It binds each statement to a theory, background, phase space, correction or
boundary class, lifecycle state, limitation, and source certificate.  This
paper has its own claim map and verifier; the live programme portal remains
Paper 98.  A claim that is not linked to evidence is left unpromoted rather
than inferred from neighboring calculations.

Three aspects of this workflow matter scientifically.  First, a failed
promotion is retained: the invalid Berger round-sphere Hodge subspace, for
example, was replaced by full $SU(2)_L\times U(1)_R$ isotypical blocks rather
than silently repaired.  Second, causal, spectral, algebraic, and quantum
results carry different dependency labels.  Third, stopping rules are fixed
before the result is known.  The counterflow candidate passed its causal gate
and then failed a later physical gate; both facts remain part of the record.

\appendix
\section{Frozen claim-boundary index}

Table~\ref{tab:claim-boundaries} is a compact crosswalk from the synthesis to
the authoritative Phase-1 ledger.  The identifiers are stable machine-facing
names, while the paper uses conventional scientific language.  ``Terminal''
means that the declared Phase-1 question has an outcome; it does not promote
the result beyond its scope tuple.

\begin{longtable}{@{}p{0.30\textwidth}p{0.25\textwidth}p{0.37\textwidth}@{}}
\caption{Frozen Phase-1 claim boundaries.}\label{tab:claim-boundaries}\\
\toprule
Ledger identifier & Claim state & Boundary preserved in this paper \\
\midrule
\endfirsthead
\toprule
Ledger identifier & Claim state & Boundary preserved in this paper \\
\midrule
\endhead
\cert{phase1.einstein_extra.structure} & Classified & Pre-residual compact
linear exact sequence and Lee--Wald form; no particle or final quotient. \\
\cert{phase1.taub_kuranishi.structure} & Scoped theorem frozen & Complete
finite exponential-polynomial second-order cone; full bounded locus open. \\
\cert{phase1.interaction.disposition} & Representative frozen; invariant
class open & Exact retained cyclic representative; physical action
trivialized through summed input order two; no branch-resolved class. \\
\cert{phase1.black_hole.companion} & Companion draft allowed & Horizon
curvature and finite-flux results in declared axial class and polar fixture;
no stability or universal endpoint theorem. \\
\cert{phase1.quantum.strict} & Local Euclidean one-loop obstruction & No
Lorentzian QME, state, particle, or unitarity inference. \\
\cert{phase1.quantum.compensator} & QME restored in changed theory at one loop
& Formal $\tau$-adic local Euclidean statement; no equivalence to strict
theory or unconditional all-loop result. \\
\cert{phase1.counterflow.causal_parent} & Causal only & Exact causal parent at
the selected Berger point; no positivity follows. \\
\cert{phase1.counterflow.physical_viability} & Terminally obstructed for the
declared family & Persistent $j=\tfrac12$ Hamiltonian--Hopf block rejects the
same-field family; no all-architecture no-go. \\
\cert{phase1.counterflow.clock_charge} & Classified & Fixed charge makes $D$
null and removes the clock; unrestricted clock leaves $D$ charged. \\
\cert{phase1.observer.disposition} & Draft allowed; no promotion & No
operational counterflow frequency ratio; not a receiver nonexistence theorem. \\
\bottomrule
\end{longtable}

\section{Conclusion}

The phase-one programme began with a common but underspecified question:
whether the additional poles of pure Weyl gravity are physical ghosts.  Its
main conceptual result is that there is no single promotion from pole to
particle.  Local solutions, reduced classical directions, nonlinear tangent
directions, and quantum states are different mathematical objects.

That distinction did not automatically rescue the theory.  Additional
compact-product directions survive local reduction, yet some fail bounded
nonlinear continuation.  Horizon regularity admits additional Schwarzschild
curvature solutions, yet a certified finite-norm fixture selects the
Einstein sector.  Strict pure Weyl gravity has a local Euclidean one-loop BV
obstruction; a compensator removes it only in a changed theory.  The first
changed model with a complete causal BV parent is itself rejected by an exact
family-persistent Hamiltonian--Hopf instability.

The most defensible conclusion is therefore neither ``the propagator already
settled the physics'' nor ``constraints removed every ghost.''  It is:
\begin{equation}
\boxed{
\begin{minipage}{0.84\textwidth}
Pole-level signs are an initial diagnostic.  Physical viability is a sequence
of reductions and existence tests.  In the laboratories completed in Phase
1, different additional sectors survive different early tests, but no
candidate reaches a positive full quantum state, and the first causally
complete changed theory fails a robust linear physical-stability test.
\end{minipage}}
\end{equation}

\begin{thebibliography}{99}

\bibitem{BenderMannheim}
C.~M.~Bender and P.~D.~Mannheim,
``No-ghost theorem for the fourth-order derivative Pais--Uhlenbeck oscillator
model,'' \emph{Phys. Rev. Lett.} \textbf{100} (2008) 110402,
\href{https://arxiv.org/abs/0706.0207}{arXiv:0706.0207}.

\bibitem{Salvio}
A.~Salvio,
``A non-perturbative and background-independent formulation of quadratic
gravity,'' \href{https://arxiv.org/abs/2404.08034}{arXiv:2404.08034}.

\bibitem{Anselmi}
D.~Anselmi,
``Fakeons, quantum gravity and the correspondence principle,''
\href{https://arxiv.org/abs/1911.10343}{arXiv:1911.10343}.

\bibitem{BeniniMusanteSchenkel}
M.~Benini, M.~Musante, and A.~Schenkel,
``Green hyperbolic complexes on Lorentzian manifolds,''
\href{https://arxiv.org/abs/2207.04069}{arXiv:2207.04069}.

\bibitem{BransonGover}
T.~P.~Branson and A.~R.~Gover,
``The conformal deformation detour complex for the obstruction tensor,''
\href{https://arxiv.org/abs/math/0605192}{arXiv:math/0605192}.

\bibitem{GoverSombergSoucek}
A.~R.~Gover, P.~Somberg, and V.~Sou\v cek,
``Yang--Mills detour complexes and conformal geometry,''
\href{https://arxiv.org/abs/math/0606401}{arXiv:math/0606401}.

\bibitem{BrillDeser}
D.~R.~Brill and S.~Deser,
``Instability of closed spaces in general relativity,''
\emph{Commun. Math. Phys.} \textbf{32} (1973) 291--304.

\bibitem{FischerMarsden}
A.~E.~Fischer and J.~E.~Marsden,
``Linearization stability of the Einstein equations,''
\emph{Bull. Am. Math. Soc.} \textbf{79} (1973) 997--1003.

\bibitem{Moncrief}
V.~Moncrief,
``Spacetime symmetries and linearization stability of the Einstein equations
I, II,'' \emph{J. Math. Phys.} \textbf{16} (1975) 493--498 and
\textbf{17} (1976) 1893--1902.

\bibitem{AMM}
J.~M.~Arms, J.~E.~Marsden, and V.~Moncrief,
``Symmetry and bifurcations of momentum mappings,''
\emph{Commun. Math. Phys.} \textbf{78} (1981) 455--478.

\bibitem{Maldacena}
J.~Maldacena,
``Einstein gravity from conformal gravity,''
\href{https://arxiv.org/abs/1105.5632}{arXiv:1105.5632}.

\bibitem{LuPangPope}
H.~L\"u, Y.~Pang, and C.~N.~Pope,
``Conformal gravity and extensions of critical gravity,''
\href{https://arxiv.org/abs/1106.4657}{arXiv:1106.4657}.

\bibitem{DeserJoungWaldron}
S.~Deser, E.~Joung, and A.~Waldron,
``Partial masslessness and conformal gravity,''
\href{https://arxiv.org/abs/1208.1307}{arXiv:1208.1307}.

\bibitem{MomenniaHendi}
M.~Momennia and S.~H.~Hendi,
``Quasinormal modes of black holes in Weyl gravity: electromagnetic and
gravitational perturbations,''
\href{https://arxiv.org/abs/1910.00428}{arXiv:1910.00428}.

\bibitem{BarnichBrandtHenneaux}
G.~Barnich, F.~Brandt, and M.~Henneaux,
``Local BRST cohomology in gauge theories,''
\emph{Phys. Rept.} \textbf{338} (2000) 439--569,
\href{https://arxiv.org/abs/hep-th/0002245}{arXiv:hep-th/0002245}.

\bibitem{FradkinTseytlin}
E.~S.~Fradkin and A.~A.~Tseytlin,
``Renormalizable asymptotically free quantum theory of gravity,''
\emph{Nucl. Phys. B} \textbf{201} (1982) 469--491.

\bibitem{Rachwal}
L.~Rachwa\l,
``Introduction to quantization of conformal gravity,''
\href{https://arxiv.org/abs/2204.13856}{arXiv:2204.13856}.

\bibitem{BaumeKerenZur}
F.~Baume and B.~Keren-Zur,
``The dilaton Wess--Zumino action in higher dimensions,''
\href{https://arxiv.org/abs/1307.0484}{arXiv:1307.0484}.

\bibitem{Dittrich}
B.~Dittrich,
``Partial and complete observables for Hamiltonian constrained systems,''
\emph{Gen. Rel. Grav.} \textbf{39} (2007) 1891--1927,
\href{https://arxiv.org/abs/gr-qc/0411013}{arXiv:gr-qc/0411013}.

\end{thebibliography}

\end{document}
